\chapter{引言}\label{chap:intro}

\section{研究背景与动机}\label{sec:bg}

家谱是研究中国前现代人口史与社会史的核心一手史料，连续记录多代人口属性\citep{zhao2001chinese}。本研究真正面向的目标域 B，是由各宗族自行编纂并流传至今的田野家谱集合：其体例严格遵循父系单线原则，嫁入家族的女性多仅以姓氏入籍，嫁出者近乎隐没；女子入谱书法因时代而异，部分时期女儿、妻、妾仅作占位姓氏出现，跨宗族婚姻与母系世系因此被系统性遮蔽\citep{yan2009nuzi,calderon2025analysing}。结果是 B 在文献层面碎裂为彼此孤立的父系树森林，且不附带可比对的婚姻真值，难以直接训练监督模型。

与 B 形成对照的是带标注训练源 A，即中国多代人口面板数据集（China Multi-Generational Panel Dataset，CMGPD-LN）\citep{campbell2011data}：基于辽宁清代户籍档案整理而成，覆盖约 151 万条人年记录、跨度 1749--1909 年，得益于户口编纂逻辑而保留了夫妻关系与母子、母女关系的完整登记。在 B 中失踪、却被 A 保留的女性，本文沿用历史档案学界术语称为母家（mujia）\citep{zhao2001chinese}。识别母家所对应的婚姻边并非简单"补缺"：婚姻形成受粮价、灾荒、家户结构、八旗归属、官员任职等多尺度因素共同塑造\citep{shiue2022marriage,batyra2023union}。传统上，史学家依靠档案爬梳\citep{duff2003list}或众包\citep{shan2025reexamining}逐对补全，难以扩展到 A 的量级；而对 B 这类无婚姻真值的目标域，更需要能在 A 上学习并迁移至 B 类场景的方法链路。本文采取的策略是：在 A 上对母系边与目标年婚姻边实施"父系镜像消融"，构造 B 真实场景的简化代理，在受控条件下评估系统的可识别性与可解释性。

\section{核心挑战}\label{sec:challenge}

近年来，得益于图卷积网络\citep{kipf2016gcn}与注意力机制\citep{vaswani2017attention}的进展，将婚姻边补全建模为异构图上的链路预测成为可行选项\citep{hu2020hgt,zhang2018seal}。但将其直接套用于家谱场景仍面临三重挑战。其一是不可解释：图神经网络通常仅输出标量分数，专家难以判断候选对是因共居、同旗还是世代深度被推荐，无从完成史学审慎\citep{calderon2025analysing}。

其二是大语言模型（Large Language Model，LLM）多步推理在大规模异构图上的局限：基于生成式智能体\citep{park2023generative}与思维链\citep{wei2022cot}的多智能体协商\citep{liang2023llm}虽具开放域建模能力，但难以原生处理百万级节点的异构图，且其推理链常被证明是事后合理化\citep{turpin2023faithfulness}。其三是史学验证的需求：任何加入家谱的边都将影响后续世代的预测，系统必须配套来源标记、协商记录与一键回滚，使每条边均附带可追溯、可推翻的证据链。

\section{研究方法概述}\label{sec:approach-overview}

针对上述挑战，本文提出名为 GeneaLink 的可视分析系统，以带标注训练源 A 上的父系镜像消融为受控代理，探索 B 类目标域的母家婚姻边重建路径。系统由三件计算骨架与一套可视分析界面构成。异构图变换器（Heterogeneous Graph Transformer，HGT）\citep{hu2020hgt}在含 4 类节点、9 类边的异构图上学习候选对排序，本文训练得到的编码器在 1880--1909 测试集上达到 Hit@10$=0.624$；子图链路预测（SEAL）\citep{zhang2018seal}对候选抽取封闭子图并施以双半径节点标注，得到 13 类与历史学语义相关的微结构图样作为可解释证据；多智能体模拟（Multi-Agent Simulation，MAS）将上述证据与档案事件注入夫妻两侧智能体，以印象生成、深入询问、质疑、对齐、复述、终评六轮协商得到对称惩罚下的最终评分。可视分析层由 6 个联动视图组成，使专家可批量接受高置信对、逐对审视边界候选、注入领域知识并回滚错误决策。

\section{论文贡献与组织}\label{sec:contribution}

本文主要贡献可归纳为四点。第一，针对母家边缺失的难题，提出"消融而不删除"的方案，在保留元数据快照的前提下对带标注训练源 A 实施父系镜像消融，使母系信号的边际贡献得以受控量化；该方案是目标域 B 真实场景的简化代理，并非等价镜像。第二，将 HGT、SEAL 与 LLM 多智能体协商耦合为端到端可解释流水线，把拓扑信号转译为智能体可援引的论据，将链路预测由黑箱评分转化为可被审视的论证过程。第三，设计由性能概览、流程视图、蜂窝画布、二分图视图、协商竞技场与规则视图组成的混合主动可视分析界面，配套来源标记、提示注入与一键回滚。第四，通过案例研究与专家评估初步验证系统的可用性与可解释性，并讨论向其他社会网络场景的迁移可能。

本文其余部分组织如下。第二章综述相关工作；第三章报告形成性研究与设计目标；第四章给出系统概述；第五章详述后端引擎，依次涵盖异构图建模、HGT 编码与婚姻评分、SEAL 子图模式提取与多智能体协商协议；第六章呈现六视图可视分析设计；第七章简述系统实现要点；第八章给出案例研究与专家评估；第九章讨论局限并展望未来工作。
